\chapter{Kac - Moody algebras and its representations}

\section{Kac - Moody algebras}
\subsection{Main definitions}

Key definitions that we need to remember:

\begin{itemize}
    \item(Generalized Cartan matrix). A matrix $A=(a_{ij})\in M_n(\bbZ)$ is said to be a Generalized Cartan matrix if it satisfies:
    \begin{equation}
        \left\{\begin{aligned}
            &a_{ii} = 2,\;\forall i\geq1\\
            &a_{ij} \leq 0,\;\forall i\neq j\\
            &a_{ij} = 0 \iff a_{ji} = 0
        \end{aligned}\right.
    \end{equation}
    \item(Realization of a matrix). If $A = (a_{ij})\in M_n(\bbZ)$ is a matrix of rank $\ell$, then a realization of $A$ is defined as a triple 
    $(\fr{h},\Pi,\check{\Pi})$, where $\Pi=\{\alpha_1,\dots,\alpha_n\}\subset\fr{h}^*$ and 
    $\check{\Pi}=\{\calpha_1,\dots,\calpha_n\}$ satisfies:
    \begin{equation}
        \left\{\begin{aligned}
            &\Pi, \check{\Pi} \;\text{ are L.I sets}\\
            &\langle \alpha_i,\calpha_j \rangle = a_{ij}
        \end{aligned}\right.
    \end{equation} 
    We claim that under this condition we have $\dim(\fr{h}) \geq 2n - \ell$. Assume that $(\fr{h},\Pi,\check{\Pi})$ is a 
    realization of $A$ and write $\dim(\fr{h})= n+m$. So, we can extend $\check{\Pi}$ to a basis of $\fr{h}$ by adding 
    L.I vectors $v_1,\dots, v_m$. Write $b_{ij} = \langle\alpha_i,v_j\rangle$, then since $\Pi$ is a linearly independent 
    set, the matrix $ M = (A\;B),\;B=(b_{ij}) \in M_{nm}(\bbC)$ must be of rank $n$. However, by reordering the indexes, we can 
    assume that $A = (\begin{smallmatrix}A_1\\A_2\end{smallmatrix}), B = (\begin{smallmatrix}B_1\\B_2 \end{smallmatrix})$, where 
    $A_1\in M_{\ell,n}, B_2\in M_{n-\ell,m}$ and $A_1$ has full rank. Then we have
    \begin{equation}
        M = \begin{pmatrix}
            A_1 & B_1 \\
            A_2 & B_2
        \end{pmatrix},\; \text{Rank}(M) = n. 
    \end{equation}
    Therefore $\text{Rank}(B_2) = n - \ell\Rightarrow m\geq n-\ell$. This proves that $\dim(\fr{h})\geq 2n - \ell$. In general 
    we will assume that a realization is minimum in the sense that $\dim(\fr{h}) = 2n - \ell$.
    \item(Existence of realizations). Given a matrix $A=(a_{ij}) \in M_n(\bbC)$, by reordering indexes we can write 
    $A = (\begin{smallmatrix} A_1\\ A_2\end{smallmatrix}), \text{Rank}(A_1) = \ell$. So the matrix 
    \begin{equation}
        C = \begin{pmatrix}
            A_1 & 0 \\
            A_2 & \text{I}_{n-\ell, n - \ell}
        \end{pmatrix}\in M_{n, 2n - \ell}
    \end{equation}
    has full rank, so its rows are linearly independent. 
    Then we let $\calpha_j = R_j(C)$, the rows of $C$ and $\alpha_i$ the first $n$ cordinate functions. This clearly 
    gives a minimum realization of $A$.
    \item(Free Lie algebras). Let $X$ be a set and $\bbC X$ the associated vector space of maps 
    $\{X\to \bbC : \text{fin. support}\}$. 
    We write $F(X) = T(\bbC X)$, which is an associative algebra. We also can view this as a Lie algbra with bracket 
    $[a,b] = ab - ba$. 
    \begin{dang}
        \begin{definition}
            The free Lie algebra $FL(X)$ generated by $X$ is the smallest Lie subalgebra of $F(X)$ containing $X$. Or equivalently, it's 
            the intersection of all Lie subalgebrs of $F(X)$ containing $X$. 
        \end{definition}
    \end{dang}
    The imporantance of this definition lies in the following proposition.
    \begin{dang}
        \begin{lemma}
            If $\theta:X\to \fr{g}$ is a map from $X$ to a Lie algebra $\fr{g}$, then there exists an unique 
            homomorphism $\phi:FL(X)\to \fr{g}$ such that $\phi|_X = \theta$.
        \end{lemma}
    \end{dang}
    Proof: By the universal property of the tensor algebra, the map $\theta : X \to \fr{g}$ induces a map of associative algebras
    $F(X)\to T(\fr{g})$. So we have an induced map $\phi:F(X)\to \mathcal{U}(\fr{g})$. We see that 
    $\phi^{-1}(\fr{g})\subset F(X)$ is a Lie subalgebra which contains $X$, so $FL(X)\subset \phi^{-1}(\fr{g})$ and we conclude that 
    $\phi: FL(X)\to \fr{g}$. It's clear that $\phi|_X = \theta$. The uniqueness follows from the fact that if $\phi^\prime$
    is another map $FL(X)\to \fr{g}$ such that $\phi^\prime|_X = \theta$, then the set of elements $a\in F(X)$ such that 
    $\phi(a)=\phi^\prime(a)$ is a Lie algebra containing $X$, so $\phi|_{FL(X)} = \phi^\prime|_{FL(X)}$.\qed
    \item(Pre Kac - Moody algebra). Let $A=(a_{ij})\in M_n(\bbC)$ and $(\fr{h},\Pi,\check{\Pi})$ be a realization of $A$. We let 
    \begin{equation}
        X = \{f_i,e_i:i=1,\dots,n\}\cup\{\fr{h}\}.
    \end{equation}
    So we have the free Lie algebra $FL(X)$. We observe that $i:\fr{h}\hookrightarrow FL(X)$ is just an inclusion of sets, 
    not a vector space monomorphism. In order to make it a vector space monomorphism, we let 
    \begin{equation}
        \fr{h} \hookrightarrow \frac{FL(X)}{\langle i(\lambda x -\mu y) - \lambda i(x) + \mu i(y)  \rangle}.
    \end{equation}
    Now we add the relations (we will identify $\fr{h}$ and $i(\fr{h})$):
    \begin{equation}
        \left\{\begin{aligned}
            &[h, e_i] = \langle \alpha_i, h\rangle e_i\\
            &[h, f_i] = -\langle\alpha_i,h\rangle f_i\\
            &[e_i,f_j] = \delta_{ij}\calpha_i\\
        \end{aligned}\right.
    \end{equation}
    for all $h\in\fr{h}$. Call all these relations by $R$, we will write 
    \begin{equation}
        \widetilde{\fr{g}}(A) = \frac{FL(X)}{R}.
    \end{equation}
    \item(A tensor algebra representation for $\widetilde{\fr{g}}(A)$). Let $V=\bigoplus_{i=1}^n\bbC v_i$, there is a 
    natural construction of a representation of $\widetilde{\fr{g}}(A)$ on $T(V)$ for each $\lambda\in\fr{h}^*$. 
    The defining relations are as follows:
    \begin{equation}
        \left\{\begin{aligned}
            &f_i(a) = v_i\otimes a,\; a\in T^s(V)\\
            &h(1) = \langle\lambda,h\rangle 1\\
            &e_i(1) = 0
        \end{aligned}\right.
    \end{equation}
    Then we use the ``creation'' operator $f_i$ to deduce the action of $h$ and $e_i$ on $T(V)$ inductively.
    \begin{center}
        \textcolor{red}{It's important to note that the defined action does not looks like a derivation ! 
        That is, we dont have in general the 
        relation:
        \begin{equation}
            X (a\otimes b) = (Xa)\otimes b + a\otimes(Xb),\;a\in T^s(V),b\in T^r(V).
        \end{equation}}
    \end{center}
    For example, we see that 
    \begin{equation}
        \begin{split}
            h(v_i) = hf_i(1) &= f_ih(1) + [h,f_i](1)\\
            &= \langle \lambda, h\rangle f_i(1) - \langle \alpha_i, h\rangle f_i(1)\\
            & = \langle \lambda - \alpha_i, h\rangle v_i. 
        \end{split} 
    \end{equation}
    Furthermore 
    \begin{equation}
        \begin{split}
            h(v_i\otimes a) &= f_ih(a) + [h,f_i](a)\\
            &=  - \langle \alpha_i, h\rangle f_i(a) + v_i\otimes h(a),\; a\in T^s(V).
        \end{split}
    \end{equation}
    In general we see that:
    \begin{equation}
        \left\{\begin{aligned}
            &e_i: T^s(V) \to T^{s-1}(V): v_j\otimes a\mapsto \delta_{ij}\calpha_j(a)+v_j\otimes e_i(a)\\
            &f_i: T^s(V) \to T^{s+1}(V): a\mapsto v_i\otimes a\\
            &h: T^s(V) \to T^s(V):v_j\otimes a \mapsto - \langle \alpha_j, h\rangle v_j\otimes a + v_j\otimes h(a),\;h\in\fr{h}.
        \end{aligned}\right.
    \end{equation}
    In addition we see that $T(V)$ is a weight representation:
    \begin{equation}
        \begin{split}
    &h \cdot v_{i_1}\otimes\cdots\otimes v_{i_k}=\langle\lambda - \sum_{j=1}^k\alpha_{i_j}, h \rangle v_{i_1}\otimes\cdots\otimes v_{i_k}\\
    &\quad\Rightarrow T(V) = \bigoplus_{k_1,\dots,k_n\geq0} V_{\lambda - \Sigma_{j=1}^nk_j\alpha_j}.        
        \end{split}
    \end{equation}

    Using this family of representations we deduce the next corollary.
    \begin{dang}
        \begin{corollary}
            We have $\widetilde{\fr{g}}(A) = \widetilde{\fr{n}}_-\oplus\fr{h}\oplus\widetilde{\fr{n}}_+$, where 
            $\widetilde{\fr{n}}_\pm$ are the subalgebras generated by the $e'$s and $f'$s respectively.
        \end{corollary}
    \end{dang}
    Proof: It's easy to see that $\widetilde{\fr{g}}(A) = \widetilde{\fr{n}}_-+\fr{h}+\widetilde{\fr{n}}_+$, now let 
    $n_- + h + n_+=0\in \widetilde{\fr{g}}(A)$ and $\lambda\in\fr{h}^*$. 
    Using the representation constructed above we see that $n_-(1) + \langle\lambda, h\rangle 1 = 0$. This last expression 
    is true for all $\lambda\in\fr{h}^*$, then $h = 0$. Now it remains to prove that 
    \begin{equation}
       \theta: \widetilde{\fr{n}}_- \to T(V): n_- \mapsto n_-(1)
    \end{equation}
    is injective. To see this we observe that the above map is just the restriction of the  $(T(V),\lambda)$
    restricted to $\widetilde{\fr{n}}_-$, so it's a Lie algebra homomorphism. Its image contains $\{v_1,\dots,v_n\}$ so 
    \begin{equation}
        \theta:\widetilde{\fr{n}}_- \to FL(V).
    \end{equation}
    However, this map has an inverse given by $v_i\mapsto f_i$. Then we conclude that $\theta$ is injective.\qed
    \item(More notation). We have the set of notations.
    \begin{equation}
        \left\{\begin{aligned}
            & Q = \sum_i\bbZ\alpha_i, Q_+ =\sum_i\bbZ_{\geq0}\alpha_i,\;Q_- = -Q_+\\
            &\fr{g}_\alpha = \{x\in \widetilde{\fr{g}}(A): [h,x]=\langle\alpha,h\rangle x\}\\
            &(\text{These implies:})\; \widetilde{\fr{g}}(A) = (\bigoplus_{\alpha\in Q_-\setminus0}\fr{g}_\alpha)\oplus \fr{h}
            \oplus (\bigoplus_{\alpha\in Q_+\setminus 0})\fr{g}_\alpha\\
            &\fr{g}_j(s) = \bigoplus_\alpha\fr{g}_\alpha,\;\alpha=\sum_ik_i\alpha_i,\;\sum_ik_is_i = j\\
            &\fr{g}_j = \fr{g}_j(1,\dots,1) = \bigoplus_{\alpha\;:\;\text{ht}(\alpha) = j}\fr{g}_\alpha
        \end{aligned}\right.
    \end{equation}
    \item(The Kac - Moody algebra). We define 
    \begin{equation}
        \fr{g}(A) = \frac{\widetilde{\fr{g}}(A)}{ I }
    \end{equation}
    where $I$ is the maximal ideal which intersects $\fr{h}$ trivially. This maximal ideal exists because the general result that 
    \begin{equation}
        V = \bigoplus_\lambda V_\lambda \Rightarrow W = \bigoplus_\lambda (W\cap V_\lambda)
    \end{equation}
    for any weight representation $V$ of a Lie algebra and a submodule $W\subset V$. Therefore, the sum of ideals 
    intersecting $\fr{h}$ trivially, is an ideal intersecting $\fr{h}$ trivially.

    More notation:
    \begin{equation}
        \left\{\begin{aligned}
            &\fr{g}_\alpha = \{x\in\fr{g}(A):[h,x]=\langle\alpha,h\rangle x,\;\forall h\in\fr{h}\}\\
            &\text{mult}(\alpha) = \dim(\fr{g}_\alpha)\\
            &\Delta_+ = \{\alpha\in Q_+: \text{mult}(\alpha) \neq 0\},\;\Delta_-=-\Delta_+\\
            &(\text{The root system:})\;\Delta = \Delta_+ \cup \Delta_-
        \end{aligned}\right.
    \end{equation}
\end{itemize}

\subsection{Invariant bilinear form}

A generalized Cartan matrix $A$ is said to be symmetrizable if there exist a decomposition
\begin{equation}
    A = DB,\; D = \text{diag}(\ep_1,\dots,\ep_n), B=(b_{ij}) \text{ symmetric}.
\end{equation}
We also impose the condition on $D$ being invertible, i.e., $\ep_i\neq0\;\forall i$.\par 

Let $\fr{g} = \fr{g}(A)$ be a Kac - Moody algebra, with $A$ being a symmetrizable generalized Cartan matrix. We define 
a bilinear form $(\cdot|\cdot):\fr{h}\times\fr{h}\to\bbC$ as follows:
\begin{equation}\label{equation : invairant bilinear form relations}
    \left\{\begin{aligned}
        &\fr{h} = \fr{h}^\pr + \fr{h}^{\pr\pr},\;\fr{h}^{\pr} = \sum_i\bbC\calpha_i\\
        &(\calpha_i,|h)=\langle\alpha_i,h\rangle\ep_i,\;\forall h\in\fr{h}\\
        &(h^{\pr}|h^{\pr\pr}) = 0 ,\;\forall h^{\pr},h^{\pr\pr}\in\fr{h}^{\pr\pr}
    \end{aligned}\right.
\end{equation} 

\begin{dang}
    \begin{theorem}
        There exist an unique symmetric bilinear form on $\fr{g}(A)$ that's invariant, nondegenerate on $\fr{h}$ and satisfying 
        relations~(\ref{equation : invairant bilinear form relations}).
    \end{theorem}
\end{dang}
Proof: If $(\cdot|\cdot)$ satisfies~(\ref{equation : invairant bilinear form relations}), then it's nondegenerate on $\fr{h}$. 
In fact, if $(\sum_ic_i\calpha_i|h)=0,\;\forall h\in\fr{h}$, 
then $\langle \sum_ic_i\ep_i\alpha_i, h\rangle= 0,\;\forall h\in\fr{h}$. Therefore
\begin{equation}
    \sum_ic_i\ep_i\alpha_i = 0\Rightarrow c_i = 0,\;\forall i.
\end{equation}
Now, by invariance we see that this form must satisfies $(x|y)=0$ whenever 
$(x,y)\in\fr{g}_\alpha\times\fr{g}_\beta,\alpha+\beta\neq0$. In fact, for any $h\in\fr{h}$, we have 
\begin{equation}
    \langle \alpha,h\rangle(x|y) = ([h,x]|y)=-(x|[h,y])=-\langle\beta,h\rangle(x|y)
\end{equation}
and $\langle\alpha+\beta,h\rangle(x|y)=0$. Since $\alpha+\beta\neq0$, we see that $(x|y)=0$. From this we also deduce that 
$(\cdot|\cdot)|_{\fr{g}_\alpha\times\fr{g}_{-\alpha}}$ is nondegenerate, otherwise $(\cdot|\cdot)$ will be degenerate, which means 
that $\Ker(\cdot|\cdot)$ would be a nonzero ideal that intersects $\fr{h}$ trivially, which is a contradiction.\par 

We also can use invariance to find a formula for $(\cdot|\cdot)$ on $\fr{g}(N) = \bigoplus_{j=-N}^N\fr{g}_j$. First, on 
$\fr{g}(1) = \sum_i\bbC f_i+\fr{h}+\sum_i\bbC e_i$, we already know that we must have 
\begin{equation}
    \left\{\begin{aligned}
        &(\fr{g}_{\pm1},\fr{g}_{\pm1}) = 0\\
        &(\fr{g}_{\pm1},\fr{h}) = 0
    \end{aligned}\right.
\end{equation}
It still remains to calculate $(e_i|f_j)$. By invariance, we must have 
\begin{equation}
    \langle\alpha_i,h\rangle(e_i|f_j)=(e_i|[f_j,h])=([e_i,f_j]|h) = \delta_{ij}(\calpha_i|h).
\end{equation}
So, for $h=\calpha_j$ we deduce the relation $(e_i|f_j)=\delta_{ij}\ep_i$. Now, assume that we have defined 
$(\cdot|\cdot)$ on $\fr{g}(N-1)$. To extend this form to $\fr{g}(N)$ we need to define it on $\fr{g}_N\times\fr{g}_{-N}$. 
So let $(x,y)\in\fr{g}_N\times\fr{g}_{-N}$ and write $x=\sum_i[u_i,v_i]$ where $\deg(u_i)+\deg(v_i)=N, u_i,v_i\in\fr{g}(N-1)$. 
Although $([u_i,v_i], y)$ is not yet defined, we can use invariance to write 
\begin{equation}
    (x|y) = \sum_i(u_i|[v_i,y])
\end{equation}
because $\deg([v_i,y])>-N$. That turns out to be a definition for the extension of the bilinear form from 
$\fr{g}(N-1)$ to $\fr{g}(N)$. \qed

Since $(\cdot|\cdot)$ is nondegenerate, we have an induced map $\nu:\fr{h}\to\fr{h}^*$. It's clear from equation~(\ref{equation : invairant bilinear form relations})
that $\nu(\calpha_i)=\ep_i\alpha_i$. Therefore
\begin{equation}\label{equation: epsilon in terms of the bilinear form}
    \begin{split}
        (\alpha_i|\alpha_i) &=\ep_i^{-2}(\calpha_i|\calpha_i)\\
        &=\ep_i^{-2}\ep_i\langle\calpha_i,\alpha_i\rangle\\
        &=\ep_i^{-1}2\\
        &\Rightarrow\ep_i = \frac{2}{(\alpha_i|\alpha_i)}.
    \end{split}
\end{equation}
So we conclude that $\calpha_i = \tfrac{2}{(\alpha_i|\alpha_i)}\nu^{-1}(\alpha_i)$ and 
$A_{ij} = 2\tfrac{(\alpha_j|\alpha_i)}{(\alpha_j|\alpha_j)}$.

\subsection{Weyl group of a Kac - Moody algebra}

A quick sketch for this subsection:
\begin{center}
    \begin{tikzcd}
        & \text{Representation theory of }\fr{sl}_2\bbC \arrow[ld]\arrow[dd]\\
     \ad(e_i),\;\ad(f_i) \text{ are nilpotent operators on } \fr{g}(A)\arrow[d]\arrow[rd, dashed] & \\
     W = \biggl\langle r_i = \exp(f_i)\exp(-e_i)\exp(f_i) \biggr\rangle \arrow[r, dashed]&   \dim(V_\lambda) = \dim(V_{r_i(\lambda)})
    \end{tikzcd}
\end{center}


\begin{dang}
    \begin{lemma}
        On a Kac - Moody algebra $\fr{g}(A)$, the operators $ad(e_i)$ and $\ad(f_i)$ are 
        locally nilpotent.
    \end{lemma}
\end{dang}
Proof: The proof include several steps:
\begin{itemize}
    \item(A little lemma). If $a\in\fr{n}_+$ is such that $[f_i,a]=0,\;\forall i$, then $a = 0$. In fact, if $a\neq0$, then 
    \begin{equation}
         \sum_{i,j\geq0} \ad(\fr{g}_1)^{i}\ad(\fr{h})^ja
    \end{equation}
    would be a nonzero ideal intersecting $\fr{h}$ trivially, which is a contradiction.

    \item(Representation theory of $\fr{sl}_2$). Recall that if $V$ is a $\fr{sl}_2\bbC=\bbC\langle Y,H,X\rangle$
    module and $v\in V$ is a vector satisfying:
    \begin{equation}
        \left\{\begin{aligned}
            &H(v) = \lambda v\\
            &X(v) = 0
        \end{aligned}\right.
    \end{equation}
    Then $\lambda = n\in\bbZ_{\geq0}$ and $XY^k(v) = k(n - k + 1)Y^{k-1}(v)$. In fact, $W = \langle Y^k(v):k\geq0 \rangle $
    is an irreducible $\fr{sl}_2\bbC-$module.
    \item(Serre relations). We have 
    \begin{equation}
        \ad(e_i)^{1-a_{ij}}e_j = 0 = \ad(f_i)^{1 - a_{ij}}f_j\;\text{for }i\neq j.
    \end{equation}
    We will shot that $\ad(f_i)^{1 - a_{ij}}f_j = 0$ by proving that it commutes with all the $e$'s. If we write 
    $v = f_j$ and think $\fr{g}(A)$ as a representation under the adjoint action of 
    $\fr{sl}_2\bbC = \bbC\langle f_i, \calpha_i, e_i\rangle$, then 
    \begin{equation}
        e_i(v) = 0,\; \calpha_i(v) = -a_{ij}v.  
    \end{equation}  
    Then we see that 
    \begin{equation}
        \begin{split}
            [e_i, \ad(f_i)^{1 - a_{ij}}f_j] &= e_i\cdot f_i^{1-a_{ij}}(v)\\ 
            &= (1 - a_{ij}) ( - a_{ij}  - (1 - a_{ij}) + 1)f_{i}^{ - a_{ij}} (v)\\
            &= 0. 
        \end{split}
    \end{equation}
    Moreover $[e_k, \ad(f_i)^{1 - a_{ij}}f_j] = 0,\forall k\neq i,j$. Finally, we have
    \begin{equation}
        \begin{split}
             [e_j, \ad(f_i)^{1 - a_{ij}}f_j] &= \ad(f_i)^{1-a_{ij}}\calpha_j\\
             &= a_{ij}\ad(f_i)^{-a_{ij}}(f_i)\\
             &= 0,\;\text{if }a_{ij} < 0. 
        \end{split}
    \end{equation}
    For $a_{ij} = 0$ we still have $[e_j, [f_i,f_j]] = [f_i, \calpha_j] = a_{ij}\calpha_j = 0$.
    \item(Conclusion). Said this, we know that for each $j$, there exist $N_j\geq1$ such that 
    $\ad(f_i)^{N_j}f_j = 0$. The same is true for the other generators of $\fr{g}(A)$. Finally, an induction on the length
    of a Lie word on $f$'s and $e$'s proves that $\ad(f_i)$ is locally nilpotent on $\fr{g}(A)$.
\end{itemize}

\begin{dang}
    \begin{definition}
        We define the Weyl group of $\fr{g}(A)$ by 
        \begin{equation}
            W = \biggl\langle \theta^{\ad}_i = e^{\ad(e_i)}e^{-\ad(f_i)}e^{\ad(e_i)} \big|_{\fr{h}} \biggr\rangle \subset GL(\fr{h}).
        \end{equation}
    \end{definition}
\end{dang}

\begin{dang}
    \begin{lemma}
        We have 
        \begin{equation}\label{equation : Weyl group action on coroots}
    \theta_i(h) = h - \langle \alpha_i, h\rangle \calpha_i.
         \end{equation}
    \end{lemma}
\end{dang}
Proof: It's clear that if $\langle\alpha_i, h\rangle = 0$, then $\exp(\ad(e_i))h= h = \exp(\ad(f_i))h$. Therefore, in this case we 
have $\theta_i(h) = h$. So it remains to show that $\theta_i(\calpha_i) = -\calpha_i$. To see this we use the 
standard representation of $\fr{sl}_2\bbC$ to write 
\begin{equation}
    e_i = \begin{pmatrix}
        0&1\\0&0
    \end{pmatrix},
    f_i = \begin{pmatrix}
        0&0\\1&0
    \end{pmatrix},
    \calpha_i = 
    \begin{pmatrix}
        1&0\\0&-1
    \end{pmatrix}.
\end{equation}
So, we see that 
\begin{equation}
    \exp(e_i) = \begin{pmatrix}
        1&1\\0&1
    \end{pmatrix},
    \exp(-f_i) = \begin{pmatrix}
        1&0\\-1&1
    \end{pmatrix}.
\end{equation}
Therefore
\begin{equation}
    \begin{split}
        \theta^{\ad}_i(\calpha_i) & = \Ad_{\exp(e_i)}\Ad_{\exp(-f_i)}\Ad_{\exp(e_i)}\calpha_i\\
        &= \begin{pmatrix}
            0&1\\-1&0
        \end{pmatrix}
        \begin{pmatrix}
            1&0\\0&-1
        \end{pmatrix}
        \begin{pmatrix}
            0&-1\\1&0
        \end{pmatrix}\\
        &=\begin{pmatrix}
            -1&0\\0&1
        \end{pmatrix}\\
        &=-\calpha_i.
    \end{split}
\end{equation}
\qed

Taking the contragradient representation of $W$ on $\fr{h}^*$ and using the 
formula~(\ref{equation : Weyl group action on coroots}), we arrive at an equivalent definition for the Weyl group.

\begin{dang}
\begin{definition}\label{definition: Weyl group on roots}
    We define the Weyl group of $\fr{g}(A)$ by the group generated by reflections defined by 
    \begin{equation}
        r_i(\lambda) = \lambda - \langle\lambda, \calpha_i \rangle\alpha_i.
    \end{equation}
\end{definition}
\end{dang}

\begin{dang}
    \begin{corollary}
        Assume that $V$ is an integrable $\fr{g}(A)-$module, i.e., $V$ is a weight representation
        \begin{equation}
            V = \bigoplus_{\lambda\in\fr{h}^*}V_\lambda
        \end{equation}
        and each $e_i,f_i$ acts on $V$ as nilpotent operators.  
        Then $\theta_i(V_\lambda) = V_{r_i(\lambda)}$, where we let
        \begin{equation}
            \theta_i = \exp(e_i)\exp(-f_i)\exp(e_i)\in \End(V).
        \end{equation} 
        In particular $\dim(V_{w(\lambda)}) = \dim(V_\lambda),\;\forall w\in W$.
    \end{corollary}
\end{dang}
Proof: First, remember some formulas. Let $X,T\in\End(V)$ be nilpotent operators on $V$. Then we have 
\begin{equation}
    \begin{split}
        \exp(T)X\exp(-T) &= \sum_{i,j\geq0}\frac{(-1)^j}{i!j!}T^iXT^j\\
        &=\sum_{m\geq0}\frac{1}{m!}\sum_{j=0}^m\binom{m}{j}(-1)^jT^{m-j}XT^j\\
        &=\sum_{m\geq0}\frac{1}{m!}(L_T - R_T)^mX\\
        &=\sum_{m\geq0}\frac{1}{m!}\ad(T)^mX\\
        &=\exp(\ad(T))X.
    \end{split}
\end{equation}

If $\theta_i = \exp(e_i)\exp(-f_i)\exp(e_i)$, then by the previous equation we see that 
\begin{equation}
    \begin{split}
        \theta_i\cdot\calpha_i\cdot\theta_i^{-1} &= \exp(\ad(e_i))\exp(\ad(-f_i))\exp(\ad(e_i))\calpha_i\\
        &=\theta_i^{\ad}(\calpha_i)\\
        &=-\calpha_i.
    \end{split}
\end{equation} 
Let $v\in V_\lambda$, then it's clear that $h\cdot\theta_i(v) = \langle \lambda,h\rangle \theta_i(v)$ because 
$h$ commute with $e_i$ and $f_i$. However, by the previous calculation we see that we also have 
$\calpha_i\theta_i(v) = -\theta_i\calpha_i(v)=-\langle \lambda,\calpha_i\rangle\theta_i(v)$. Therefore
\begin{equation}
    \begin{split}
        (h - \frac12\langle &\lambda, h\rangle\calpha_i)\cdot\theta_i(v) =
        \langle \lambda, h - \frac12\langle \lambda, h\rangle\calpha_i\rangle \theta_i(v)\Rightarrow\\
        h\cdot\theta_i(v) &= \langle\lambda, h - \langle\alpha_i,h\rangle\calpha_i\rangle\theta_i(v)\\
        & = \langle\lambda,\theta_i^{\ad}(h)\rangle \theta_i(v)\\
        &= \langle r_i(\lambda),h\rangle \theta_i(v).
    \end{split}
\end{equation}
\qed

\begin{dang}
    \begin{lemma}
        The invariant bilinear form defined on $\fr{g}(A)$ is preserved under the action of the Weyl group.
    \end{lemma}
\end{dang}
Proof: Let $a,b,x\in\fr{g}(A)$, then we deduce (from invariance):
\begin{equation}
    (\ad(x)^{i}a|b) = (-1)^i(a| \ad(x)^{i}b).
\end{equation}
Therefore 
\begin{equation}
    (\exp(ad(x))a|\exp(\ad(x))b) = (a|\exp(-\ad(x))\exp(\ad(x))b) = (a|b).
\end{equation}
So we see that $(\theta_i^{\ad}(a)|\theta_i^{\ad}(b))=(a|b)$.\qed

\subsection{Affine Lie algebras}
\begin{itemize}
    \item(Affine Cartan matrix). An affine GCM is a GCM satisfying:
    \begin{equation}
        \begin{aligned}
           &(1)\;\text{Corank}(A) = 1\\
           &(2)\;Au = 0 \text{ for some } u>0\\
           &(3)\; Au\geq0\Rightarrow Au = 0
        \end{aligned}
    \end{equation}
    For example, lets classify $2\times 2$ generalized Cartan matrices of affine type. If 
    $A=(\begin{smallmatrix}2&a\\b&2\end{smallmatrix})$ is affine, then by the first condition it has rank 1, so 
    $\det(A)=4-ab=0$ and since $a,b\in\bbZ_{\leq0}$ we see that $a\in\{-1,-2,-4\}$. We see two possibilities (take into acount the transposes):
    \begin{equation}
        \begin{pmatrix}
            2&-2\\-2&2
        \end{pmatrix}\text{ and }
        \begin{pmatrix}
            2&-1\\-4&2
        \end{pmatrix}.
    \end{equation}    
    We still need to see that the other two conditions are satisfied. For the second one we take $u=(1,1)$ and $u=(1,2)$ for 
    the two matrices respectively. The third condition is clearly satisfies for both matrices.
    \item(Realization by loop algebras extensions). If $\fr{g}$ is a finite dimensional Lie algebra with 
    an invariant bilinear form $(\cdot,\cdot)$, we have affinization given by 
    \begin{equation}
        \widehat{g} = \fr{g}[t^\pm] +\bbC K,\;[xt^m,yt^n]=[x,y]t^{m+n} + m\delta_{m,-n}(x,y)K.
    \end{equation}
    The Kac - Moody affinization, is an extension of $\widehat{\fr{g}}$ by a derivation. We have 
    \begin{equation}
        \widetilde{g} = \widehat{g} +\bbC d,\; [d, K] = 0,[d,xt^m] = mxt^m.
    \end{equation}
    In the Kac - Moody affinization, the vector space $\widetilde{\fr{h}} = \fr{h} +\bbC K + \bbC d$ plays the role of Cartan 
    subalgebra. We will show that $\widetilde{g}$ is an affine Kac moody algebra.
    \item(Explicty realization for affine Cartan matrices). Let $\fr{g}$ be a simple finite dimensional Lie algebra and choose 
    a Cartan subalgebra $\fr{h}\subset \fr{g}$ and a set of positive roots. Then we have a triangular decomposition 
    \begin{equation}
        \fr{g} = \bigoplus_{\alpha\in\Delta_+}\fr{g}_{-\alpha}+\fr{h}+\bigoplus_{\alpha\in\Delta_+}\fr{g}_\alpha.
    \end{equation}
    Let $\Pi =\{\alpha_1,\dots,\alpha_\ell\}$ be the set of simple roots (so $\text{Rank}(\fr{g})=\ell$) and choose 
    $(E_i,F_i)\in\fr{g}_{\alpha_i}\times\fr{g}_{-\alpha_i}$ in such a way that $\langle F_i, H_i, E_i \rangle$, where 
    $H_i=[E_i,F_i]$, is a $\fr{sl}_2-$triple. So $H_i = \calpha_i$, where we define in general 
    \begin{equation}\label{equation: coroots definition}
        \begin{aligned}
            &\nu :\fr{h}\overset{\cong}{\to}\fr{h}^*:H\mapsto K(H,\;\cdot\;),\\
            &\alpha^\vee = \frac{2}{(\alpha,\alpha)}\nu^{-1}(\alpha),\;\alpha\in\Delta.
        \end{aligned}
    \end{equation} 
    Let $\theta\in\Delta_+$ be the (unique) maximal root of maximal heith: Remember that $\Delta_+\subset\bbZ_{\geq0}\Pi$ and 
    $\text{ht}(\sum_ik_i\alpha_i)=\sum_ik_i,\; \sum_ik_i\alpha_i\in\Delta_+$. For example, in $\fr{g}=\fr{sl}_{n+1}\bbC$
    we have a triangular decomposition given by 
    \begin{equation}
        \left\{\begin{aligned}
            &\fr{g} = \bigoplus_{i>j}\fr{g}_{L_i - L_j} + \fr{h} + \bigoplus_{i < j}\fr{g}_{L_i - L_j},\\
            &\fr{h} = \text{ traceless diagonal matrices},\\
            &\fr{g}_{L_i - L_j} = \bbC E_{ij},\;i\neq j\\
            &\Pi = \{\alpha_1,\dots,\alpha_n\},\;\alpha_i = L_i - L_{i+1}\\
            &\theta = \sum_{i=1}^n \alpha_i = L_1 - L_{n+1}\\
            &K(\cdot,\cdot) = 2n\Tr(\dot,\cdot).
        \end{aligned}\right.
    \end{equation}
    In particular, since $\langle E_{n+1,1}, H_{1,n+1}, E_{1,n+1} \rangle,\; H_{1,n+1}= E_{11} - E_{n+1,n+1}$, we have 
    $\theta^\vee = H_{1,n+1}$.\par 

    So given a simple finite dimensional Lie algebra, and $\widetilde{\fr{g}}$ its Kac - Moody affinization, we define 
    \begin{equation}
        \begin{aligned}
            &\Pi^\vee =\{\underset{\calpha_0}{\underbrace{\tfrac{2}{(\theta,\theta)}K - \theta^\vee\otimes1}},
            \calpha_i: i =1,\dots,\ell\}\text{ where }\calpha_i=H_i\otimes1\;(i>0),\\
            &\Pi=\{\underset{\alpha_0}{\underbrace{\delta - \theta}},\alpha_i:i=1,\dots,\ell\}.
        \end{aligned}
    \end{equation}
    Here we observe that we extend $\fr{h}^* \hookrightarrow \widetilde{\fr{h}}^*$ by putting $\alpha|_{\bbC K+\bbC d} = 0$ for 
    $\alpha\in\fr{h}^*$. Furthermore we define $\delta\in\widetilde{\fr{h}}^*$ by $\delta(d) = 1,\delta|_{\fr{h}+\bbC K}=0$.\par 

    Define $A = (\langle\alpha_i,\calpha_j\rangle)_{i,j=0}^\ell$, we claim that it's an affine GCM (\TODO proof). For example, 
    for $\fr{g} = \fr{sl}_{n+1}\bbC$ we have the induced affine GCM:
    \begin{equation}
        \begin{pmatrix}
            2&-1&0&\cdots&0&-1\\
            -1&2&-1&\cdots&0&0\\
            0&-1&2&\cdots&0&0\\
            \vdots&\vdots&\vdots&\cdots&\vdots&\vdots\\
            0&0&0&\cdots&2&-1\\
            -1&0&0&\cdots&-1&2
        \end{pmatrix}
    \end{equation}
    
    Before we prove that $\widetilde{g}$ is the affine Kac - Moody algebra associated to the affine GCM $A$, we need to define 
    a set of generators for it. We let 
    \begin{equation}
        \left\{\begin{aligned}
            &e_0 = E_{-\theta}\otimes t,\;f_0=E_{\theta}\otimes t^{-1}\\
            &e_i = E_i\otimes 1,\;f_i = F_i\otimes 1\\
        \end{aligned}\right.
    \end{equation}
    To understand this convention, we can look to the root system of $\widetilde{\fr{sl}_2}$:
    \begin{center}
        \begin{tikzcd}
            \vdots & \vdots& \vdots \\
           \bullet_{F\otimes t^2} & \bullet_{H\otimes t^2}& \bullet_{E\otimes t^2} \\
           \textcolor{red}{\bullet_{F\otimes t}} & \bullet_{H\otimes t}& \bullet_{E\otimes t} \\
           \bullet_{F\otimes 1} & \bullet_{H\otimes1}\arrow[lu]\arrow[r]& \textcolor{red}{\bullet_{E\otimes1}} \\
           \bullet_{F\otimes t^{-1}} & \bullet_{H\otimes t^{-1}}& \bullet_{E\otimes t^{-1}} \\
           \vdots & \vdots& \vdots \\
        \end{tikzcd}
    \end{center}
    We see in general that 
\begin{equation}\label{equation: root system of an affine Lie algebra}
    \begin{aligned}
        &\widetilde{\Delta} = \{\alpha + j\delta: \alpha\in\Delta,j\in\bbZ\}\cup\{j\delta:j\neq0\}\\
        &\widetilde{g}_{\alpha + j\delta} = \bbC E_\alpha\otimes t^j, \widetilde{g}_{j\delta} = \fr{h}t^j.
    \end{aligned}
\end{equation}
\item(Main theorem).
Before we prove that $\widetilde{\fr{g}}$ is an Affine Kac - Moody algebra, we need a lemma.
\begin{dang}
    \begin{lemma}
        Let $\fr{g}$ and suppose that:
        \begin{itemize}
            \item $\fr{h}\subset\fr{g}$ is a commutative subalgebra
            \item $\fr{g} =\bbC\langle e_i, f_i: (i=1,\dots,n), \fr{h}\rangle$ 
            \item There exist $\Pi = \{\alpha_1,\dots,\alpha_n\}\subset\fr{h}^*$ and 
            $\Pi^\vee=\{\calpha_1,\dots,\calpha_n\}\subset\fr{h}$ (both L.I sets) such that 
            \begin{equation}
                \begin{aligned}
                    &[h,e_i]=\langle\alpha_i,h\rangle e_i\\
                    &[h,f_i]=-\langle\alpha_i,h\rangle f_i\\
                    &[e_i,f_i] = \delta_{ij}\calpha_i
                \end{aligned}
            \end{equation}
            for all $h\in\fr{h}$.
            \item Any nonzero ideal of $\fr{g}$ intersects $\fr{h}$ trivially
        \end{itemize}
        Then $(\fr{g},\fr{h},\Pi,\Pi^\vee)$ is a realization of $A=(\langle\alpha_i,\calpha_j\rangle)_{i,j=1}^n$.
    \end{lemma}
\end{dang}
Proof: Since $\fr{g}$ is generated by $\fr{h}$ and $e_i,f_i\;(i=1,\dots,n)$ we have a surjective map 
\begin{equation}
\varphi:\widetilde{\fr{g}}(A)\twoheadrightarrow \fr{g}.
\end{equation}
Now $I=\Ker(\varphi)\subset \widetilde{\fr{g}}(A)$ is an ideal. Since $\varphi|_{\fr{h}} = \mathds{1}_{\fr{h}}$ by construction, 
we see that $I\cap\fr{h} = (0)$. Moreover, if $J\subset \widetilde{\fr{g}}(A)$ intersecting $\fr{h}$ trivially, then 
$\varphi(J)\subset \fr{g}$ is an ideal intersecting $\fr{h}$ trivially, so $\varphi(J) = 0$ and $J\subset\Ker(\varphi)$. 
Therefore $\Ker(\varphi)$ is the maximal ideal of $\widetilde{\fr{g}}(A)$ intersecting $\fr{h}$ trivially and we have an 
induced isomorphism $\fr{g}(A) \overset{\cong}{\rightarrow}\fr{g}$.\qed

\begin{dang}
    \begin{theorem}
        Let $\fr{g}$ be a finite dimensional simple Lie algebra. Choose a Cartan subalgebra 
        $\fr{h}\subset\fr{g}$ and the induced triangular decomposition
        \begin{equation}
            \fr{g} = \fr{h}\oplus\bigoplus_{\alpha\in\Delta}\fr{g}_\alpha.
        \end{equation} 
        For a choice of positive roots we have a set of simple roots $\{\alpha_1,\dots,\alpha_\ell\}$ and the associated
        set of coroots $\{H_1,\dots,H_\ell\}$ defined by equation~(\ref{equation: coroots definition}). 
        Let $\theta$ be the (unique) root of maximal heith, then we define 
        \begin{equation}
            \alpha_0 = \delta - \theta,\;\calpha_0 = \frac{2}{(\theta,\theta)}K - \theta^\vee\otimes1
        \end{equation}
        where $\delta\in\widetilde{\fr{h}}^*,\widetilde{\fr{h}}=\fr{h}+\bbC K+\bbC d, \delta|_{\fr{h}+\bbC K}=0,\delta(d)=1$. 
        Furthermore, let $\Pi=\{\alpha_0,\dots,\alpha_\ell\}$ and $\Pi^\vee=\{\calpha_0,\dots,\calpha_\ell\}$, we claim that 
        the quadruple $(\widetilde{\fr{g}}, \widetilde{\fr{h}},\Pi,\Pi^\vee)$ is a realization of the affine GCM 
        $(\langle\alpha_i,\calpha_j\rangle)_{i,j=0}^{\ell}$.
    \end{theorem}
\end{dang}
Proof: This proof have three parts:
\begin{itemize}
    \item($\widetilde{\fr{g}}$ is generated by the $e_i,f_i\;(i=0,\dots,\ell)$). Remember that we define 
    \begin{equation}
        e_0 = E_{-\theta}\otimes t, f_0=E_{\theta}\otimes t^{-1}, e_i = E_i\otimes 1, f_i=F_i\otimes 1 (i=1,\dots,\ell)
    \end{equation}
    where $E_i = E_{\alpha_i}, F_i = E_{-\alpha_i}$ are such that $\fr{sl}_2\cong\bbC\langle F_i,H_i,E_i\rangle$. Since 
    $\theta$ is of maximal heith and $\fr{g}$ is simple, we have 
    $\fr{g}=\mathcal{U}(\fr{g})E_{-\theta}=\mathcal{U}(\fr{n}_+)E_{-\theta}$. Now define $\widetilde{g}_1$ as the Lie 
    algebra generated by $\widetilde{\fr{h}},e_i,f_i\;(i=0,\dots,\ell)$. Since $\fr{g}$ is simple, it's generated by 
    the $E_i,F_i$'s and we have $\fr{g}\subset\widetilde{\fr{g}}$. 
    Furthermore we see $\mathcal{U}(\fr{n}_+)e_0=\fr{g}\otimes t\in\widetilde{\fr{g}}_1$, and an inudction show that 
    $\fr{g}\otimes t^{j}\in \widetilde{\fr{g}}_1,\;\forall j\geq0$. In a similar way we see that 
    $\fr{g}\otimes t^{j}\in\widetilde{\fr{g}}_1,\;j<0$. So we conclude that $\widetilde{\fr{g}}_1=\widetilde{\fr{g}}$.

    \item($\widetilde{\fr{g}}$ generators satisfies all the relations). 
    \begin{equation}
     \begin{split}
         &[e_0,f_0] = -\theta^\vee\otimes 1 + \frac{2}{(\theta,\theta)}K = \calpha_0\\
         &[e_0, f_i] = [E_{-\theta}\otimes t, F_i\otimes1]=[E_{-\theta},F_i]\otimes t = 0\;(\text{ht}(\theta)\text{ is minimum})\\
         &[e_0, h] =[E_{-\theta},h]\otimes t = \theta(h)e_0 =-\langle\delta-\theta,h\rangle e_0\;(\forall h\in\fr{h})\\
         &[e_0,d] = -e_0 = -\langle\delta-\theta,d\rangle e_0\\
         &[e_0,K] = 0 = -\langle\delta-\theta, K\rangle e_0\\
         &\text{Etc}\dots
     \end{split}
    \end{equation}
    \item(Any nonzero ideal intersects $\widetilde{h}$ trivially). Let $I\subset\widetilde{\fr{g}}$ be a nonzero ideal. Since 
    $I=\bigoplus_{\alpha+j\delta\neq0}I\cap\widetilde{\fr{g}}_{\alpha+j\delta}$ we know that there exist 
    $x\otimes t^j\in I\cap \fr{g}_{\alpha+j\delta}\setminus0$ 
    because of equation~(\ref{equation: root system of an affine Lie algebra}). 
    Since the Killing form is nondegenerate, there exist $y\in \fr{g}_{-\alpha}$ such that 
    $(x|y)\neq0$. Therefore 
    \begin{equation}
        [y\otimes t^{-j},x\otimes t^j] = [x,y]\otimes1 -j(x|y)K\in\widetilde{\fr{h}}\Rightarrow j=0
    \end{equation}
    and $[x,y] = (x|y)\nu^{-1}(\alpha) = 0\Rightarrow \alpha = 0$. This constradicts the hypothesis $\alpha+j\delta\neq0$.
    \qed
\end{itemize} 
\item(Invariant bilinear form).
By the previous theorem we know that $\widetilde{g}$ is a Kac - Moody algebra, so it has an unique invariant bilinear $(\cdot|\cdot)$ form satisfying 
equation~(\ref{equation : invairant bilinear form relations}). On Cartan $\widetilde{\fr{h}}$ we must have 
\begin{equation}\label{equation: invariant form on affine Cartan (relations)}
    \begin{aligned}
        &(d|d) = 0\\
        &(d|\calpha_i) = \langle\alpha_i,d\rangle=0\\
        &(\tfrac{2}{(\theta,\theta)}K-\theta^\vee\otimes1|\tfrac{2}{(\theta,\theta)}K-\theta^\vee\otimes1) =2\\
        &(\tfrac{2}{(\theta,\theta)}K-\theta^\vee\otimes1| d) = \ep_0\langle \delta-\theta, d\rangle =\ep_0\\
        &(\tfrac{2}{(\theta,\theta)}K-\theta^\vee\otimes1| \calpha_i) = -\ep_0\langle\theta,H_i\rangle\;(i>0)
    \end{aligned}
\end{equation}
However, we know that $\fr{h}=\text{Span}(\{H_1,\dots,H_\ell\})$ and $\langle\alpha_i,\bbC K+\bbC d\rangle = 0,\;i>1$. Therefore, 
we have $(\fr{h}\otimes1|\bbC K+\bbC d)=0$. From the previous equation we see that 
\begin{equation}
    \begin{aligned}
        &(\tfrac{2}{(\theta,\theta)}K|\tfrac{2}{(\theta,\theta)}K) = 0\Rightarrow (K|K) = 0\\
        &(\tfrac{2}{(\theta,\theta)}K|d) = \ep_0\Rightarrow (K|d) = \ep_0\frac{(\theta,\theta)}{2}
    \end{aligned}
\end{equation}
\begin{remark}
    To avoid circular arguments, observe that $(\cdot|\cdot)|_{\fr{g}}$ is an invariant bilinear form. So choose a normalization 
    such that $(\theta|\theta)=(\theta,\theta)$, where 
    \begin{equation}
        (\cdot,\cdot) = \frac{1}{2h^\vee}K(\cdot,\cdot)\;(h^\vee \text{ is the dual coxeter number}).
    \end{equation}
\end{remark}

By equation~(\ref{equation: invariant form on affine Cartan (relations)}), we also have
\begin{equation}
   \begin{split}
    &\frac{2}{(\theta|\theta)}(K|\calpha_i) - (\theta^\vee\otimes1|\calpha_i) = -\ep_0\langle\theta,H_i\rangle\\
    &\iff\frac{2\ep_i}{(\theta|\theta)}\underset{=0}{\underbrace{\langle\alpha_i,K\rangle}} - (\theta^\vee\otimes1|\calpha_i) = -\ep_0\langle\theta, H_i\rangle\\
    &\iff\ep_0\langle\theta,H_i\rangle = (\theta^\vee\otimes1|H_i\otimes1)\\
    &\iff \ep_0\langle\theta,H_i\rangle=\frac{2}{(\theta|\theta)}(\nu^{-1}(\theta)|H_i)\\
    &\iff \ep_0\langle\theta,H_i\rangle = \frac{2}{(\theta|\theta)}\langle\theta,H_i\rangle\\
    &\iff\ep_0 = \frac{2}{(\theta|\theta)}.
   \end{split}
\end{equation}
From the previous remark we have $(\theta|\theta)=(\theta,\theta)$ and $\ep_0=1$. Maybe we can avoid all this calculation 
by simply choosing $\ep_0=1$, because the symmetrized form of a GCM can be modified by multiplying the symmetric part 
by a scalar matrix.\par 


Then we see that $(K|d)=1$ and by definition 
we have $\nu(K) = \delta$. We also define $\Lambda_0 = \nu(d)$, so we have 
\begin{equation}
    \begin{aligned}
        &(\delta|\delta) = (d|d) = 0\\
        &(\Lambda_0|\Lambda_0) = (K|K) = 0\\
        &(\Lambda_0|\delta) = (K|d) = 1
    \end{aligned}
\end{equation}

Now we determine $(\cdot|\cdot)$ on $\widetilde{g}$. It's sufficient to determine the pairing where it's different from zero. So we 
calculate $(x\otimes t^j|y\otimes t^{-j}), (x,y)\in\fr{g}_\alpha\times\fr{g}_{-\alpha}$:
\begin{equation}
    \begin{split}
    \langle\alpha,H_i\rangle(x|y)&=([H_i,x]|y)\\
    &=(H_i\otimes1|[x,y]\otimes1)\\
    &=([H_i,x]\otimes t^j|y\otimes t^{-j})\\
    &=\langle\alpha,H_i\rangle (x\otimes t^j|y\otimes t^{-j})
    \end{split}
\end{equation}
If $\alpha\neq0$, then $\alpha(H_i)\neq0$ for some $i$ and we are done. So we conclude that 
\begin{equation}
    (x\otimes t^m|y\otimes t^n) = \delta_{m,-n}(x|y).
\end{equation}
This is compatible with the general expression:
\begin{equation}
    (x\otimes P|y\otimes Q) = \text{Res}_t(t^{-1}PQ)(x|y).
\end{equation}
\item (The Weyl group). First lets compute the Weyl group of $\widetilde{sl}_2$. In this case we have 
\begin{center}
    \begin{tabular}{|c|}
        \hline 
        $\fr{sl}_2 = \bbC F+\bbC H + \bbC E$\\
        \hline 
        $\alpha = \theta: H\mapsto 2$\\
        \hline
        $\theta^\vee = H$\\
        \hline
        $\Pi^\vee = \{K-H\otimes1,H\otimes1\}$\\
        \hline
        $\Pi=\{\delta-\alpha, \alpha\}$\\
        \hline
    \end{tabular}
\end{center}
Remember the definition of 
the Weyl group acting on roots~(\ref{definition: Weyl group on roots}). 
\begin{equation}
    \begin{split}
        &r_0=\left\{\begin{aligned}
        &\alpha\mapsto \alpha-\langle\alpha,K-H\otimes1\rangle(\delta-\alpha) = 2\delta-\alpha\\
        &\Lambda_0\mapsto \Lambda_0-\langle\Lambda_0,K-H\otimes1\rangle(\delta-\alpha)=\Lambda_0-\delta+\alpha\\
        &\delta\mapsto \delta - \langle\delta,K-H\otimes1\rangle(\delta-\alpha)=\delta\\
    \end{aligned}\right.\\
    &\Rightarrow r_0(m\alpha+k\Lambda_0+f\delta)=(k-m)\alpha+k\Lambda_0+(\delta-k+2m)\delta
    \end{split}
\end{equation}
We see that in the basis $\{\alpha,\Lambda_0,\delta\}$ we have 
\begin{equation}
    r_1 = \text{Diag}(-1,1,1)\text{ and } 
    r_0 = \begin{pmatrix}
        -1&1&0\\0&1&0\\2&-1&1
    \end{pmatrix} 
\end{equation}
The Weyl group of $\widetilde{sl}_2$ is $W=\langle r_0,r_1\rangle$.\par 

Now lets describe the structure of Weyl groups of affine Lie algebras in general. We will write 
$\widetilde{\lambda}=\lambda + n\delta+k\Lambda_0\in\widetilde{\fr{h}}^*$, where $\lambda=\text{Proj}(\widetilde{\lambda})\in\fr{h}^*$. 
From definition~(\ref{definition: Weyl group on roots}), we see that 
\begin{equation}
   \left\{ \begin{aligned}
        &r_0(\widetilde{\lambda})=\widetilde{\lambda} -\frac{2}{(\theta|\theta)}(k-(\theta|\lambda))(\delta-\theta)
        \;(\text{because }\langle\delta-\theta,\widetilde{\lambda}\rangle= k-(\theta|\lambda))\\
        &r_i(\widetilde{\lambda})=\widetilde{\lambda} - 2\frac{(\lambda|\alpha_i)}{(\alpha_i|\alpha_i)}\alpha_i
    \end{aligned}\right.
\end{equation}
becasue $(\widetilde{\lambda}|h)=(\lambda|h),\;\forall h\in\fr{h}$. Furthermore, we \textit{define}
\begin{equation}
    \begin{split}
        r_\theta(\widetilde{\lambda}) &= \widetilde{\lambda}-2\frac{(\widetilde{\lambda}|\theta)}{(\theta|\theta)}\theta\\
        &=\widetilde{\lambda}-2\frac{(\lambda|\theta)}{(\theta|\theta)}\theta.
    \end{split}
\end{equation}
Now lets compute 
\begin{equation}
    \begin{split}
        r_0 r_\theta(\widetilde{\lambda})&= r_0(\widetilde{\lambda}-2\frac{(\lambda|\theta)}{(\theta|\theta)}\theta)\\
        &=\widetilde{\lambda} - 2\frac{(\lambda|\theta)}{(\theta|\theta)}\theta-\frac{2}{(\theta|\theta)}
        (k-(\theta|\lambda)+2\frac{(\theta|\lambda)(\theta|\theta)}{(\theta|\theta)})(\delta-\theta)\\
        &=\widetilde{\lambda}+k\frac{2\theta}{(\theta|\theta)}-(\langle\lambda,\theta^\vee\rangle+k\frac{\|\theta^\vee\|^2}{2})\delta \\
        &=\widetilde{\lambda}+k\nu(\theta^\vee)-(\langle\lambda,\theta^\vee\rangle+k\frac{\|\theta^\vee\|^2}{2})\delta.
    \end{split}
\end{equation}
This motivates the following notation:
\begin{notation}
    Let $M=\bbZ W\nu(\theta^\vee)$, we define the \textit{translation}
    \begin{equation}
        t_\alpha(\widetilde{\lambda})=\widetilde{\lambda}+k\alpha-((\lambda|\alpha)+k\frac{\|\alpha\|^2}{2})\delta,\;\alpha\in M.
    \end{equation}
\end{notation}
Observe that if we write $\alpha=\nu(\theta^\vee)$, we have $\langle\lambda,\theta^\vee\rangle = (\nu^{-1}(\lambda)|\nu^{-1}(\alpha))=(\lambda|\alpha)$
and $\|\alpha\|^2=\|\theta^\vee\|^2$, so our notation is consistent $(t_{\nu(\theta^\vee)}=t_\alpha)$.
\begin{dang}
    \begin{lemma}
        For $w\in W$ and $\alpha,\beta\in M$ we have 
        \begin{itemize}
            \item[(1)] $wt_\alpha w^{-1}=t_{w(\alpha)}$
            \item[(2)] $t_{\alpha}t_\beta=t_{\alpha+\beta}$
        \end{itemize}
    \end{lemma}
\end{dang}
Proof: For the first identity we observe that since $w\in W$ (the Weyl group of $\fr{g}$), we have 
$w(\widetilde{\lambda})=w(\lambda)+n\delta+k\Lambda_0$. Therefore
\begin{equation}
    \begin{split}
        wt_\alpha w^{-1}(\widetilde{\lambda})&=wt_\alpha(w^{-1}(\lambda)+n\delta+k\Lambda_0)\\
        &=w(w^{-1}(\lambda)+n\delta+k\Lambda_0+k\alpha-((w^{-1}(\alpha)|\alpha)+k\frac{\|\alpha\|^2}{2})\delta\\
        &=\lambda+n\delta+k\Lambda_0+kw(\alpha)-((\alpha|w(\alpha))+k\frac{\|w(\alpha)\|^2}{2})\delta\\
        &=t_{w(\alpha)}(\widetilde{\lambda}).
    \end{split}
\end{equation}
For the second one we compute:
\begin{equation}
    \begin{split}
        t_\alpha t_\beta(\widetilde{\lambda})&
        =t_\alpha(\widetilde{\lambda}+k\beta-((\lambda|\beta)+k\frac{\|\beta\|^2}{2})\delta)\\
        &=\widetilde{\lambda}+k\beta -((\lambda|\beta)
        +k\frac{\|\beta\|^2}{2})\delta+k\alpha - ((\lambda+k\beta|\alpha)+k\frac{\|\alpha\|^2}{2})\delta \\
        &=t_{\alpha+\beta}(\widetilde{\lambda}).
    \end{split}
\end{equation}

\begin{dang}
    \begin{theorem}
        The set $T=\{t_\alpha:\alpha\in M\},\; M=\bbZ W\nu(\theta^\vee)$ is a normal subgroup of 
        the Weyl group $\widetilde{W}$ of $\widetilde{\fr{g}}$. Furthermore we have 
        $\widetilde{W}=TW,\; T\cap W = \{1\}$. In one world, we have the equality:
        \begin{equation}
            \widetilde{W} = W\rtimes T.
        \end{equation} 
    \end{theorem}
\end{dang}
Proof: First observe that $t_{\nu(\theta^\vee)}=r_0r_\theta \in\widetilde{W}$ and since 
$t_{w\nu(\theta^\vee)}=wt_{\nu(\theta^\vee)}w^{-1}\in\widetilde{W}$, we see that $T\subset \widetilde{W}$. Furthermore, 
the identity $wt_\alpha w^{-1}=t_{w(\alpha)},\;\alpha\in M,w\in W$ implies that $W\subset N(T)$ and since 
$\widetilde{W}=\langle r_\theta, W\rangle=TW$, we see that $T$ is normal in $\widetilde{W}$. Finally, the quality 
$W\cap T=\{1\}$ follows because $n\in\bbZ\mapsto t_{n\alpha}$ is injective. So, since $W$ is finite, we can not have 
$t_\alpha\in W,\;\alpha\neq0$.\qed


\begin{example}
    The last theorem provide us a nice description of the Weyl group of $\widetilde{\fr{sl}}_2$. In this case we 
    have $M=\bbZ\alpha$ and 
    \begin{equation}
        \widetilde{W} = \{t_{m\alpha}\}_{m\in\bbZ}\cup\{t_{m\alpha}\circ r\}_{m\in\bbZ},\;r(\alpha)=-\alpha.
    \end{equation}
\end{example}
\end{itemize}

\subsection{The Kac character formula}
\begin{remark}
    Let $M$ be a weight module of a Kac - Moody algebra, then we write $\lambda>\mu$ if 
    \begin{equation}
        \lambda - \mu \in \sum_i\bbZ_{\geq0}\alpha_i 
    \end{equation}
    where $\Pi=\{\alpha_1,\dots,\alpha_n\}$.
\end{remark}

\begin{itemize}
    \item(Verma module). Let $\fr{g}=\fr{g}(A)$ be a Kac - Moody algebra and $\lambda\in\fr{h}^*$, we have a $\fr{h}+\fr{n}_+-$module 
    $\bbC v_\lambda$ defined by $\fr{n}_+v_\lambda=0$ and $h\cdot v_\lambda = \langle\lambda,h\rangle v_\lambda$. We let 
    \begin{equation}
        M(\lambda) = \mathcal{U}(\fr{g})\otimes_{\mathcal{U}(\fr{h}+\fr{n}_+)}\bbC v_\lambda.
    \end{equation}
    The Verma module is a weight representation. In fact, we have 
    \begin{equation}
        M(\lambda) = \bigoplus_{k_1,\dots,k_n\geq0}M(\lambda)_{\lambda-\Sigma_ik_i\alpha_i}.
    \end{equation}
    Moreover 
    \begin{center}
        $\dim(M(\lambda)_{\mu})=$ `` number of ways of expressing $\lambda - \mu$ as a sum of positive roots''
     \end{center}
    For example, let $\fr{g} = \fr{sl}_3$. We have positive roots $\{\alpha_1,\alpha_2,\alpha_3\}$ where 
    $\alpha_3 = \alpha_1+\alpha_2$. So 
    \begin{equation}
        M(\lambda)_{2\alpha_1 + \alpha_2} = \bbC F_1^2F_2 v_\lambda + \bbC F_1F_3 v_\lambda 
    \end{equation}
    So $\dim M(\lambda)_{2\alpha_1 + \alpha_2} = 2$, which is the number of ways of expressing $2\alpha_1+\alpha_2$ as a sum 
    of elements in $\{\alpha_1,\alpha_2,\alpha_3\}$.\par 
    Then we see that, in general $\dim M(\lambda)_\mu<\infty$. Moreover, each weight is bounded by $\lambda$, so in particular, 
    this means that $\dim(\mathcal{U}(\fr{n}_+)v) <\infty$ for any $v\in M(\lambda)_\mu$.\par 
    Since $M(\lambda)$ is a weight representation, we can define the maximal ideal $J(\lambda)$ 
    intersecting $\bbC v_\lambda$ trivially. The quocient $L(\lambda) = M(\lambda)/J(\lambda)$ is an irreducible module, because 
    any proper ideal of $M(\lambda)$ must intersects $\bbC v_\lambda$, so it must be contained in $J(\lambda)$. In general, each 
    highest weight module $M$ with highest weight vector $v_\lambda^\prime$ is a quotient of $M(\lambda)$. To see this, we need to 
    verify that $M(\lambda)\twoheadrightarrow M: u v_\lambda\mapsto u v_\lambda^\prime$ is well defined. (\TODO proof).

\begin{remark}
    We see that $L(0)$ is the trivial module. In fact, it's clear that $\mathcal{U}(\fr{g})\fr{g}_{-1}v_0\subset M(0)$ is 
    a proper submodule of codimension 1. In fact, for each $j$ we have $E_iF_jv_0 = 0,\;\forall i$. In particular 
    \begin{equation}
        \ch(L(0)) = 1\cdot e^0 =1.
    \end{equation}
\end{remark}



    \item(Category $\mathcal{O}$)
    The Category $\mathcal{O}$ of a Kac - Moody algebra is the full subcategory of $\fr{g}-$modules whose objects $M$ satisfies:
    \begin{itemize}
        \item[(1)] $M$ is finitely generated as a $\mathcal{U}(\fr{g})-$module
        \item[(2)] $M$ is a weight module 
        \item[(3)] The action of $\fr{n}_+$ is locally nilpotent, which is equivalent to 
        $\dim(\mathcal{U}(\fr{n}_+)v) <\infty,\;\forall v\in M$.
    \end{itemize}
    The third conditions is equivalent to $(3^\prime)$ which says that there exist a finite collection of weiths 
    $\{\lambda_1,\dots,\lambda_r\}$ which are upper bounds for the weigths of $M$, i.e., $M_\mu\neq0\Rightarrow \mu<\lambda_i$ 
    for some $i$. In fact, suppose that $(3^\prime)$ is satisfied and let $v\in M_\mu\setminus0$. Then 
    \begin{equation}
        \#\{(k_1,\dots,k_n)\in\bbZ_{\geq0}^n\text{ such that }M_{\mu +\sum_jk_j\alpha_j}\neq0\}<\infty
    \end{equation}
    because $M_{\mu +\sum_jk_j\alpha_j}\neq0\Rightarrow \sum_jk_j\alpha_j < \lambda_i-\mu$ for some $i$. 
    (\TODO proof of the converse).\par 

    It's clear that the Verma modules are objects in the Category $\mathcal{O}$. Furthermore, we have 
    $V,U\in\text{Obj}(\mathcal{O})\Rightarrow V\oplus U,V\otimes U\in\text{Obj}(\mathcal{O})$ and if 
    $U\subset V$ we also have $V/U\in\text{Obj}(\mathcal{O})$. In particular we see that highest weight modules 
    are ojects in this category.\par 
    Let $M\in\text{Obj}(\mathcal{O})$, we define a formal series 
    \begin{equation}
        \ch(M) = \sum_{\mu\in\fr{h}^*}\dim(M_\mu)e^\mu.
    \end{equation}
    \begin{example}
        Consider $\fr{g}=\fr{sl}_2$ and let $\lambda=k\oomega\in\fr{h}^*,\oomega(H)=1$. We have 
        $M(\lambda)_{\lambda - 2m\oomega} = \bbC F^m v_\lambda$ and the fundamental equation:
        \begin{equation}\label{equation: fundamental equation}
            EF^mv_\lambda = m(k-m+1)F^{m-1}v_\lambda.
        \end{equation} 
        Then we see that if $k>0$, then
        \begin{equation}
          M(\lambda) =  M(k\oomega) \supset M(-(k+2)\oomega) \supset 0
        \end{equation}
        has the property that each factor is irreducible. Let $x=e^{k\oomega}$, then 
        \begin{equation}
            \begin{split}
                &\ch(M(\lambda)) = x^k + x^{k-2} +x^{k-4}+\cdots = \frac{x^k}{1-x^{-2}}\\
                &\ch(M(-(k+2)\oomega)) = \frac{x^{-k-2}}{1-x^{-2}}
            \end{split}
        \end{equation}
        Therefore 
        \begin{equation}
            \ch(L(\lambda)) = \frac{x^k - x^{-k-2}}{1-x^{-2}} = x^{k} + x^{k-2} + \cdots + x^{-k}.
        \end{equation}
    \end{example}
    
    \begin{dang}
        \begin{lemma}
        The character of a Verma module is given as follows:
        \begin{equation}
        \ch(M(\lambda)) = \frac{e^\lambda}{\prod_{\alpha\in\Delta_+}(1-e^{-\alpha})^{m_\alpha}}.
        \end{equation}
        \end{lemma}
    \end{dang}
    Proof: For each $\alpha$, write a basis $\fr{g}_\alpha=\langle e_\alpha^{(i)}\rangle$. We know that 
    $M(\lambda)_\mu\neq0\iff \lambda-\mu\in Q_+$. In such a case $\dim(M(\lambda)_\mu)$ is equal to the number 
    of ways to write $\lambda-\mu$ as the sum of positive roots. This corresponds to basis elements of 
    $M(\lambda)_\mu$ given by  
    \begin{equation}
        e_{-\alpha_{i_1}}^{r_1}\cdots e_{-\alpha_{i_k}}^{r_k}v_\lambda,\; \sum_{j=1}^k r_j\alpha_{i_j} = \lambda-\mu.
    \end{equation}
    So, to count the number of these sequences we just need to take the coefficient of $\lambda-\mu$ in the expansion
    \begin{equation}
        \prod_{\alpha\in\Delta_+}(1+e^{-\alpha}+e^{-2\alpha}+\cdots)^{m_\alpha}.
    \end{equation}
    \qed
    \item(A filtration).
    \begin{dang}
        \begin{definition}
            Let $M\in\text{Obj}(\mathcal{O})$ and $\mu\in\fr{h}^*$, a filtration of $M$ with respect to $\mu$ is a sequence 
            \begin{equation}\label{equation: filtration}
                M=M_0\supset M_1\supset\cdots M_t = 0
            \end{equation}
            such that for each $i\geq1$, we have one of the two excludent possibilities:
            \begin{itemize}
                \item[(1)] $(M_{i-1}/M_i)_\lambda = 0,\forall \lambda > \mu$
                \item[(2)] $(M_{i-1}/M_i)_\mu\cong L(\mu)$
            \end{itemize}
        \end{definition}
    \end{dang}

    \begin{example}
        For $\fr{g}=\fr{sl}_2$ and $k>0$ we have 
        \begin{equation}
            M(\lambda) =  M(k\oomega) \supset M(-(k+2)\oomega) \supset 0.
        \end{equation}
        where $M(k\oomega)/M(-(k+2)\oomega)\cong L(k\oomega)$ and $M(-(k+2)\oomega)_{\mu} = 0\,\forall\mu>k\oomega$.
    \end{example}

    \begin{dang}
        \begin{lemma}
            (Existence of filtrations). For each module $M\in\text{Obj}(\mathcal{O})$ and $\lambda\in\fr{h}^*$ we claim 
            that such filtration~(\ref{equation: filtration}) exists. In this case, the number of ocurrence of 
            $L(\lambda)$ in the factors is independent of the filtration. We will write $[M:L(\lambda)]$ to the number 
            of ocurrences $L(\lambda)$ as a factor in any filtration for $\lambda\in\fr{h}^*$. 
        \end{lemma}
    \end{dang}
    (\TODO Proof:)

    \begin{dang}
        \begin{lemma}
            Let $M\in\text{Obj}(\mathcal{O})$ and $\lambda\in\fr{h}^*$, we have
            \begin{equation}
                \dim(M_\mu) = \sum_{\lambda\in\fr{h}^*}[M:L(\lambda)]\dim(L(\lambda)_\mu).
            \end{equation}
        \end{lemma}
    \end{dang}
    Proof: Take a filtration satisfying the condition of the definition~(\ref{equation: filtration}). So 
    \begin{equation}
        \begin{split}
            \dim(M_\mu) &=\sum_{i\geq1}\dim(M_{i-1})_\mu-\dim(M_i)_\mu \\
            &=\sum_{i\geq1}\dim(M_{i-1}/M_i)_\mu\\
            &=\sum_{\lambda>\mu}[M:L(\lambda)]\dim(L(\lambda)_\mu)\\
            &=\sum_{\lambda\in\fr{h}^*}[M:L(\lambda)]\dim(L(\lambda)_\mu)\;(\text{because }L(\lambda)_\mu=0,\;\mu
            \overset{\neq}{>}\lambda)
        \end{split}
    \end{equation}

    \item(Integrable modules).
    \begin{dang}
        \begin{definition}
            A $\fr{g}-$module $V$ is called integrable if $e_i,f_i\;(i=1,\dots,n)$ acts on $V$ as integrable operators.
        \end{definition}
    \end{dang}
    \begin{notation}
    \begin{equation}
        \begin{aligned}
        &X = \{\lambda\in\fr{h}^*: \langle\lambda,\calpha_i\rangle\in\bbZ\;\forall i\}\;(\text{integral weights})\\
        &X^+=\{\lambda\in \fr{h}^*: \langle \lambda,\calpha_i\rangle\in\bbZ_{\geq0}\;\forall i\}\;(\text{integral, dominant weights})
        \end{aligned}
    \end{equation}
    \end{notation}

    \begin{dang}
        \begin{lemma}
            $L(\lambda)$ is integrable if and only if $\lambda$ is dominant and integral.
        \end{lemma}
    \end{dang}
    Proof: The direction $(\Rightarrow)$ follows from the representation theory of $\fr{sl}_2$. Let $v_\lambda$ be a 
    highest weight vector in $L(\lambda)$. Then, since $L(\lambda)$ is integrable, $f_i^rv_\lambda=0$ for some $r>0$. It follows that 
    $\text{Span}(f_i^jv_\lambda:j\geq0)$ is a finite dimensional $\fr{sl}_2\cong\bbC\langle f_i,\calpha_i,e_i\rangle$ module.
    It highest weight is $\lambda$, so $\langle\lambda,\calpha_i\rangle\in\bbZ_{\geq0}$ as a consequence of the representation theory 
    of $\fr{sl}_2\bbC$. Now lets show the direction $(\Leftarrow)$, so assume $\lambda\in X^+$. The fundamental equation~(\ref{equation: fundamental equation})
    shows that 
    \begin{equation}
        e_if_i^{m+1}v_\lambda = 0,\; m = \lambda(\calpha_i) + 1\in\bbZ_{>0}.
    \end{equation} 
    Furthermore $e_jf_i^{m+1}v_\lambda=0,j\neq i$ while 
    $hf_i^{m+1}v_\lambda = \langle \lambda - (m+1)\alpha_i,h\rangle f_i^{m+1}v_\lambda$. Therefore 
    $\mathcal{U}(\fr{g})f_i^{m+1}v_\lambda \subset L(\lambda)$ is a proper submodule, so it should be zero. Therefore 
    $f_i^{m+1}v_\lambda = 0$. Now let $u\in\mathcal{U}(\fr{g})$, the equation
    \begin{equation}
        f_i^Nuv_\lambda = (\ad(f_i)+R_{f_i})^Nu= \sum_{j=0}^N\binom{N}{j}\ad(f_i)^{N-j}uf_i^jv_\lambda = 0  
    \end{equation}
    if $N\gg0$ as a consequence of $\fr{g}$ being an integrable module. To finnish the proof we also observe that 
    $e_i$ is locally nilpotent because $L(\lambda)\in\text{Obj}(\mathcal{O})$.\qed  

    \item(Generalized Casimir).
    Let $\rho\in\fr{h}^*$ be such that $\langle\rho,\calpha_i\rangle=1,\;\forall i$. Also take dual basis 
    $(u_i,u^i)\in\fr{h}\times\fr{h}, (e_\alpha^{(i)},e_{-\alpha}^{(i)})
    \in\fr{g}_{\alpha}\times\fr{g}_{-\alpha},\;\forall\alpha\in\Delta_+$. We define the generalized Casimir operator
    \begin{equation}
        \Omega = 2\nu^{-1}(\rho) +\sum_iu_iu^i +\sum_{\alpha\in\Delta_+}\sum_ie_{-\alpha}^{(i)}e_{\alpha}^{(i)}
    \end{equation}
    It's well defined as an endomorphism of an integrable module $M$. Furthermore it commutes with $\mathcal{U}(\fr{g})$, so 
    on $L(\lambda)$ it acts as a constant by the Schur's lemma. This constant can be easily obtained by applyin $\Omega$
    on the highest weight vector $v_\lambda$:
    \begin{equation}
        \begin{split}
            \Omega v_\lambda &= 2\langle\lambda,\nu^{-1}(\rho)\rangle 
            + \sum_i\langle\lambda,u_i\rangle\langle\lambda,u^i\rangle\\
            &= 2(\rho|\lambda)+\sum_i(\nu^{-1}(\lambda)|u_i)(\nu^{-1}(\lambda)|u^i)\\
            &=2(\rho|\lambda)+(\nu^{-1}(\lambda)|\nu^{-1}(\lambda))\\
            &=((\lambda+\rho|\lambda+\rho)-(\rho|\rho))v_\lambda
        \end{split}
    \end{equation}
    \begin{example}
        For $\widetilde{\fr{g}}$ affine, $\widetilde{\rho} = h^\vee \Lambda_0 + \rho$ works, where $h^\vee$ is the dual coxeter number 
        of $\fr{g}$ and $\rho=\tfrac12\sum_{\alpha\in\Delta_+}\alpha$ is the Weyl vector. (\TODO proof).
     \end{example}
    \item(Kac character formula).
    \item(Jacobi triple product).
\end{itemize}

